\section{模型的改进与推广}

\subsection{模型的改进}

若能获得多年销量和价格数据，可使用滚动时间窗估计各作物的边际分布，并通过样本外回测替代当前的区间分布假设。问题三可进一步引入作物交叉价格弹性或采购订单数据，对关系边的方向和强度进行统计估计。对于管理便利，可在最小面积占比之外加入地块间作业距离、机械切换成本和劳动力峰值约束，使面积表更接近实际执行条件。

求解方面，可采用滚动时域作为大规模实例的备用方法，但必须保留跨年豆类覆盖和重茬边界。若加入更多随机变量，可使用情景削减或Benders分解控制规模，并继续以完整硬约束核验和同任务基线作为验收条件。

\subsection{模型的推广}

本文框架可用于其他多地块、多品种和多周期的农业配置问题。更换地块适宜性、产销参数和轮作集合后，可扩展到设施农业、订单农业或区域种植结构调整。样本外评价部分也可用于比较保险、最低收购价和合同销售等政策情景，但必须重新确认概率假设与结论范围。
